% Fuente: Auxiliar 3, «Convexidad — Soluciones de una inecuación cuadrática».
\begin{itemize}
    \item[(a)] \textbf{(Teorema)} 
Demuestre que un conjunto es convexo si y solo si su intersección con una línea arbitraria 
\[
\{ \hat{x} + t v \mid t \in \mathbb{R} \}
\]
es convexa.
    \item[(b)] Sea \( C \subseteq \mathbb{R}^n \) el conjunto solución de una desigualdad cuadrática:

\[
C = \left\{ x \in \mathbb{R}^n \;\middle|\; x^T A x + b^T x + c \leq 0 \right\},
\]

donde \( A \in \mathbb{S}^n \), \( b \in \mathbb{R}^n \), y \( c \in \mathbb{R} \).

Muestre que \( C \) es convexo si \( A \succeq 0 \).
    
    \textbf{HINT}: \[
\text{Una matriz } A \in \mathbb{R}^{n \times n} \text{ es semidefinida positiva si y solo si:}
\]
\[
v^T A v \geq 0 \quad \text{para todo } v \in \mathbb{R}^n.
\]
\end{itemize}
